% Minimal statistical-testing insertions for the thesis LaTeX source.
% These blocks are intended to be pasted only into Chapter 4 and the
% corresponding limitation paragraph in Chapter 5.

% ---------------------------------------------------------------------
% Insert after the existing text in Section 4.3 Evaluation Metrics,
% immediately before \section{Experimental Protocol}.
% ---------------------------------------------------------------------
\subsection{Statistical Uncertainty and Effect Size Reporting}

To avoid interpreting single metric values as exact estimates, the main rate-based
metrics are reported together with statistical uncertainty. For hallucination rate,
non-hallucination rate, accuracy, grounded rate and abstention rate, 95\% Wilson
binomial confidence intervals are used. Wilson intervals are preferable to the
normal approximation for small evaluation sets and for rates close to 0 or 1.

For a rate \( \hat{p}=k/n \), where \(k\) is the number of positive cases and \(n\)
is the number of evaluated examples, the Wilson interval is computed as:

\begin{equation}
\frac{\hat{p} + \frac{z^2}{2n} \pm z
\sqrt{\frac{\hat{p}(1-\hat{p})}{n} + \frac{z^2}{4n^2}}}
{1+\frac{z^2}{n}},
\end{equation}

where \(z=1.96\) for a 95\% confidence interval. In addition to confidence
intervals, hallucination reductions are reported as absolute risk reduction
\((p_{\text{baseline}}-p_{\text{system}})\) and relative reduction
\((p_{\text{baseline}}-p_{\text{system}})/p_{\text{baseline}}\).

For answer F1, which is not a binary metric, paired bootstrap resampling is the
appropriate test when per-example scores are available. The examples are sampled
with replacement, the mean F1 difference between two systems is recomputed for
each bootstrap sample, and the 2.5th and 97.5th percentiles form a 95\% confidence
interval for the F1 difference. Since the reported tables mainly summarize aggregate
F1 values, the present analysis emphasizes Wilson intervals and effect sizes for
rate-based metrics, while repeated runs and full paired bootstrap testing are left as
future validation.

% ---------------------------------------------------------------------
% Insert after Table 4.1 in Section 4.5 Results on QReCC.
% Assumes the QReCC comparison table uses 50 examples.
% ---------------------------------------------------------------------
Statistical uncertainty was estimated for the QReCC hallucination rates using
95\% Wilson confidence intervals. The intervals were: Vanilla RAG 0.58
\([0.442, 0.706]\), SelfRefineAgentRAG 0.62 \([0.482, 0.741]\),
ReflexionAgentRAG 0.26 \([0.159, 0.396]\), Proposed with memory 0.28
\([0.175, 0.417]\), and Proposed without memory 0.18 \([0.098, 0.308]\).
Compared with Vanilla RAG, the proposed no-memory configuration reduced the
hallucination rate by 0.40 absolute points, corresponding to a relative reduction of
69.0\%. The proposed with-memory configuration reduced hallucination by 0.30
absolute points, corresponding to a relative reduction of 51.7\%. These effect sizes
suggest that the proposed verification architecture provides a meaningful reduction
in unsupported generation, although the confidence intervals also show that the
sample size is still limited.

The F1 values should be interpreted as aggregate answer-quality indicators rather
than statistically decisive claims. The proposed configurations improved F1 from
0.3341 for Vanilla RAG to 0.5238 with memory and 0.5199 without memory, but
a paired bootstrap over per-example F1 scores would be required to quantify the
uncertainty of these differences.

% ---------------------------------------------------------------------
% Insert after Table 4.2 in Section 4.6 Results on RAGTruth,
% before the explanatory paragraph beginning "The results of RAGTruth..."
% Assumes the RAGTruth comparison table uses 50 examples.
% ---------------------------------------------------------------------
For RAGTruth, 95\% Wilson confidence intervals were also computed for the
hallucination rates. The intervals were: VanillaRAG 0.30 \([0.191, 0.438]\),
SelfRefineAgentRAG 0.24 \([0.143, 0.374]\), ReflexionAgentRAG 0.24
\([0.143, 0.374]\), Proposed with memory 0.00 \([0.000, 0.071]\), and
Proposed without memory 0.00 \([0.000, 0.071]\). Relative to VanillaRAG, the
proposed configurations reduced the observed hallucination rate by 0.30 absolute
points. Although no hallucinations were detected for the proposed configurations
in this run, the confidence interval does not imply that the true hallucination
probability is exactly zero. Rather, it indicates that, under this sample size, the
upper 95\% confidence bound remains approximately 0.071.

% Optional replacement for the sentence "The results of RAGTruth strongly support..."
% Use this softer sentence in the paragraph after the RAGTruth statistics:
% The results of RAGTruth provide evidence in support of the proposed architecture.

% ---------------------------------------------------------------------
% Insert after Table 4.3 in Section 4.7 Results on the Legal Corpus.
% ---------------------------------------------------------------------
For the main 96-sample temporal legal setting, Wilson confidence intervals were
computed for the principal rate metrics. Accuracy was 0.8958 with a 95\% CI of
\([0.819, 0.942]\), grounded rate was 0.9583 with a 95\% CI of \([0.898, 0.984]\),
hallucination rate was 0.0104 with a 95\% CI of \([0.002, 0.057]\), and abstention
rate was 0.0312 with a 95\% CI of \([0.011, 0.088]\). These intervals show that
the observed hallucination rate is low, but the small number of hallucination cases
means that the estimate should still be interpreted with uncertainty.

For the 20-sample extractive no-memory setting, the hallucination rate of 0.0000
has a Wilson upper bound of 0.161. This is an important caution for small legal
subsets: observing no hallucinations in 20 examples does not prove that the true
hallucination rate is zero. It only shows that no hallucinations were detected in
that limited sample.

% ---------------------------------------------------------------------
% Insert after Table 4.4 in Section 4.8 Ablation Study.
% Assumes the ablation comparison uses 50 examples per condition.
% ---------------------------------------------------------------------
The ablation rates were also interpreted with Wilson confidence intervals. In the
strict setting, hallucination remained 0.10 both without and with memory, with a
95\% CI of \([0.043, 0.214]\). In the balanced setting, hallucination increased from
0.10 \([0.043, 0.214]\) to 0.12 \([0.056, 0.238]\), an absolute increase of 0.02.
In the relaxed setting, hallucination increased from 0.10 \([0.043, 0.214]\) to
0.16 \([0.083, 0.285]\), an absolute increase of 0.06 and a relative increase of
60\%. In the safe skeptic setting, hallucination decreased from 0.08
\([0.032, 0.188]\) to 0.04 \([0.011, 0.135]\), an absolute reduction of 0.04 and
a relative reduction of 50\%.

These confidence intervals overlap, so the ablation should not be read as a
definitive significance claim. The stronger conclusion is directional and
mechanistic: relaxed memory reuse can improve lexical overlap while increasing
unsupported content, whereas skeptic-filtered memory can reduce the observed
hallucination rate at the cost of a small decrease in F1. A paired bootstrap over
per-example F1 scores would be needed to test whether the F1 changes are stable.

% ---------------------------------------------------------------------
% Add this row inside Table 4.8 Error analysis and limitations,
% after the "Single-run evaluation" row.
% ---------------------------------------------------------------------
Statistical uncertainty & The study reports Wilson confidence intervals and effect
sizes for the main rate metrics. However, the sample sizes are still limited and most
experiments are single-run evaluations, so the results should be validated with
repeated runs and paired bootstrap tests for F1. \\

% ---------------------------------------------------------------------
% Replace the Chapter 5.3 single-run limitation paragraph with this version.
% ---------------------------------------------------------------------
Another limitation is that most experiments were performed as single-run
evaluations for each configuration. Although confidence intervals and effect sizes
were reported for the main rate-based metrics, single-run evaluation does not
measure run-to-run variation caused by stochastic generation. Future experiments
should repeat each configuration across several runs or random seeds and report
the mean and standard deviation for hallucination rate, answer F1, exact match,
groundedness, abstention rate and memory-hit rate. Paired bootstrap confidence
intervals should also be reported for F1 differences when per-example scores are
available.
